\documentclass[11pt]{article}

\usepackage[margin=1in]{geometry}
\usepackage{amsmath, amssymb}
\usepackage{booktabs}
\usepackage{enumitem}
\usepackage{xcolor}

% --- Notation shortcuts ----------------------------------------------------
\newcommand{\I}{\mathcal{I}}
\newcommand{\B}{\mathcal{B}}
\newcommand{\Svc}{\mathcal{S}}
\newcommand{\Evt}{\mathcal{E}}
\newcommand{\Bi}{\mathcal{B}_i}
\newcommand{\bigM}{M}
\newcommand{\ein}{e^{\mathrm{in}}}
\newcommand{\eout}{e^{\mathrm{out}}}

\title{Discrete Berth Allocation with a Shared Navigation Channel,\\
Vessel--Berth Compatibility, and Priority Berth Reservations}
\author{portsDFL --- Port of San Antonio}
\date{\today}

\begin{document}
\maketitle

\begin{abstract}
We schedule one week of vessel calls at the Port of San Antonio. Each vessel is
assigned to a \emph{compatible} berth and a start time; a single
\emph{navigation channel} (a shared one-vessel-at-a-time waterway) must be
transited once to enter and once to exit; and a subset of priority
\emph{service} vessels hold fixed berth-time \emph{reservations} that are kept
exclusively for them. The model is solved for a chosen operational objective:
total \emph{waiting time} or total \emph{berth idle time}. Objectives are
unweighted (vessel priority weights are deferred). The service time $\tau$ that
drives the schedule of the non-reserved vessels is supplied by an upstream
predictor.
\end{abstract}

\section{Sets, parameters, and variables}
\begin{table}[h]
\centering
\small
\begin{tabular}{@{}ll@{}}
\toprule
Symbol & Meaning \\
\midrule
$\I=\{1,\dots,N\}$ & vessels in the planning week \\
$\B=\{1,\dots,B\}$ & berths \\
$\Bi\subseteq\B$ & berths \emph{compatible} with vessel $i$ \\
$\Svc\subseteq\I$ & priority ``service'' vessels (fixed reservations) \\
$\Evt=\{(i,\mathrm{in}),(i,\mathrm{out}):i\in\I\}$ & the $2N$ channel transit events \\
\midrule
$a_i\ge 0$ & arrival time of vessel $i$ \\
$\tau_i>0$ & berth-occupation time: \emph{predicted} for $i\notin\Svc$, fixed at $D_r$ for $r\in\Svc$ \\
$b_r\in\B$ & reserved berth of service vessel $r\in\Svc$ \quad(fixed input) \\
$R_r\ge 0$ & reserved start time of service vessel $r$ \quad(fixed input) \\
$D_r>0$ & reserved (committed) berth-occupation duration of $r$ \quad(fixed input) \\
$c\ge 0$ & channel transit time (enter or exit) \\
$\bigM$ & big-$M$ constant ($\ge$ horizon $+$ max berth workload) \\
\midrule
$x_{ib}\in\{0,1\}$ & $1$ if vessel $i$ uses berth $b$ \quad(only for $b\in\Bi$) \\
$s_i\ge a_i$ & berth start time of vessel $i$ \\
$z_{ijb}\in\{0,1\}$ & $1$ if $i$ precedes $j$ at berth $b$ \\
$\ein_i\ge a_i$ & inbound (entry) channel-transit start \\
$\eout_i\ge 0$ & outbound (exit) channel-transit start \\
$y_{ef}\in\{0,1\}$ & $1$ if transit event $e$ precedes event $f$ \\
$\mathit{LC}_b,\ \mathit{FS}_b\ge 0$ & last completion / first start at berth $b$ (idle objective) \\
\bottomrule
\end{tabular}
\caption{Notation. Departure of vessel $i$ is $\eout_i+c$; transit event
$(i,k)$ occupies the channel over $[t_{i}^{k},\,t_{i}^{k}+c]$ with
$t_i^{\mathrm{in}}=\ein_i$ and $t_i^{\mathrm{out}}=\eout_i$.}
\end{table}

\section{Objective (selectable)}
Service vessels are pinned by their reservations (\S below), so the optimization
acts on the remaining non-service vessels $\I\setminus\Svc$. The model is solved
for one of two unweighted operational objectives:
\begin{align}
\text{(Waiting)} \quad & \min\ \sum_{i\in\I\setminus\Svc}\bigl(\ein_i - a_i\bigr),
   \label{eq:obj-wait}\\[2pt]
\text{(Idle)} \quad & \min\ \sum_{b\in\B}\bigl(\mathit{LC}_b - \mathit{FS}_b\bigr).
   \label{eq:obj-idle}
\end{align}
\emph{Waiting time}~\eqref{eq:obj-wait} is the total time the non-service vessels
wait at anchorage before they may enter the channel; the service duration
cancels, so only the controllable pre-berth delay is penalized. \emph{Berth idle
time}~\eqref{eq:obj-idle} is summed over all berths, so a berth's occupied span
includes any reserved blocks on it; since the total occupation $\sum_i\tau_i$ is
a constant, minimizing idle time is equivalent to minimizing the summed
per-berth occupied span $\sum_b(\mathit{LC}_b-\mathit{FS}_b)$, and the true idle
KPI is this span minus $\sum_i\tau_i$.

\section{Constraints}
\paragraph{Assignment and arrival.}
\begin{align}
& \sum_{b\in\Bi} x_{ib} = 1 && \forall i\in\I \label{eq:assign}\\
& \ein_i \ge a_i,\qquad s_i \ge a_i && \forall i\in\I\setminus\Svc \label{eq:arrival}
\end{align}
Constraint~\eqref{eq:assign} assigns each vessel to exactly one \emph{compatible}
berth; the variables $x_{ib}$ exist only for $b\in\Bi$, so incompatible
assignments are structurally excluded (for a service vessel $r$ we set
$\mathcal{B}_r=\{b_r\}$, so this constraint pins it to its reserved berth).
Constraint~\eqref{eq:arrival} is the dynamic arrival release: a non-service
vessel may neither enter the channel nor begin berth service before it has
arrived ($\ein_i\ge a_i$ and $s_i\ge a_i$). Service vessels are exempt---their
timing is fixed by the reservation~\eqref{eq:reserve}, so a reserved block is
held even if the vessel arrives earlier or later than $R_r$.

\paragraph{Navigation channel.}
\begin{align}
& s_i \ge \ein_i + c && \forall i\in\I \label{eq:ch-in}\\
& \eout_i \ge s_i + \tau_i && \forall i\in\I \label{eq:ch-out}\\
& t_f \ge t_e + c - \bigM\,(1-y_{ef}) && \forall\, e,f\in\Evt,\ e<f,\ v(e)\ne v(f) \label{eq:ch-fwd}\\
& t_e \ge t_f + c - \bigM\,y_{ef} && \forall\, e,f\in\Evt,\ e<f,\ v(e)\ne v(f) \label{eq:ch-bwd}
\end{align}
A vessel may berth only after its inbound transit finishes~\eqref{eq:ch-in}, and
may begin its outbound transit only after service completes~\eqref{eq:ch-out}.
Constraints~\eqref{eq:ch-fwd}--\eqref{eq:ch-bwd} are the channel
mutual-exclusion disjunction: for every pair of transit events of
\emph{different} vessels ($v(e)$ is event $e$'s vessel), one transit must clear
the channel before the other starts---so the channel holds exactly one vessel at
a time, in either direction. The same-vessel pair
$(i,\mathrm{in}),(i,\mathrm{out})$ is already ordered
by~\eqref{eq:ch-in}--\eqref{eq:ch-out}. Setting $c=0$ collapses the transit
constraints, recovering the berth-only model.

\paragraph{Berth sequencing (no overlap).}
\begin{align}
& z_{ijb} + z_{jib} \le 1 && \forall b,\ i<j:\ b\in\Bi\cap\B_j \label{eq:one}\\
& z_{ijb} + z_{jib} \ge x_{ib} + x_{jb} - 1 && \forall b,\ i<j:\ b\in\Bi\cap\B_j \label{eq:order}\\
& s_j \ge s_i + \tau_i - \bigM\,(1 - z_{ijb}) && \forall (i,j,b),\ i\ne j:\ b\in\Bi\cap\B_j \label{eq:prec}
\end{align}
Constraint~\eqref{eq:one} permits at most one ordering direction for a pair at
berth $b$. Constraint~\eqref{eq:order} forces an ordering whenever both vessels
are assigned to $b$: its right-hand side equals $1$ only when $x_{ib}=x_{jb}=1$,
and is non-positive (vacuous) otherwise. The big-$M$ precedence~\eqref{eq:prec}
then keeps the later vessel from starting until the earlier one finishes, so
their berth occupations cannot overlap. Precedence triples and $z$ variables are
created only for pairs that can share a berth ($b\in\Bi\cap\B_j$).

\paragraph{Priority berth reservations.}
\begin{align}
& s_r = R_r, \qquad \tau_r = D_r && \forall r\in\Svc \label{eq:reserve}
\end{align}
Each service vessel $r$ carries a reservation $(b_r,R_r,D_r)$ given as fixed
input. Together with $\mathcal{B}_r=\{b_r\}$ in~\eqref{eq:assign},
constraint~\eqref{eq:reserve} pins $r$ to its reserved berth, start time, and
occupation duration, so $r$ holds the interval $[R_r,\,R_r+D_r]$ on berth $b_r$.
Exclusivity needs no extra constraint: the no-overlap
sequencing~\eqref{eq:one}--\eqref{eq:prec}, applied between $r$ and any other
vessel eligible for berth $b_r$, keeps that interval empty of all other
vessels---the slot is \emph{left empty}, with no backfilling, regardless of
whether $r$ arrives before or after $R_r$. The reservation blocks only the
berth: service vessels still transit the shared channel like any other vessel
(\eqref{eq:ch-in}--\eqref{eq:ch-bwd}).

\paragraph{Berth idle-time linearization} (only for objective~\eqref{eq:obj-idle}).
\begin{align}
& \mathit{LC}_b \ge s_i + \tau_i - \bigM\,(1-x_{ib}) && \forall i,\ b\in\Bi \label{eq:idle-lc}\\
& \mathit{FS}_b \le s_i + \bigM\,(1-x_{ib}) && \forall i,\ b\in\Bi \label{eq:idle-fs}\\
& \mathit{FS}_b \le \mathit{LC}_b && \forall b\in\B \label{eq:idle-span}
\end{align}
Under minimization, $\mathit{LC}_b$ is driven down to the latest completion and
$\mathit{FS}_b$ up to the earliest start among vessels assigned to berth $b$;
the big-$M$ switches each bound off for vessels not at $b$. Constraint
\eqref{eq:idle-span} caps an empty berth's span at zero.

\section{Implementation pointers}
The MILP is \texttt{DiscreteBAP} in
\texttt{optimizers/src/bap\_optim/discrete\_bap.py}, with constructor argument
\texttt{objective\,$\in$\,\{"waiting","idle"\}}. The instance descriptor
\texttt{BAPInstance} (in \texttt{bap\_optim/instance.py}) carries
\texttt{arrivals}, \texttt{berth\_compat}, \texttt{channel\_time}, and the
service reservations as fixed $(b_r,R_r,D_r)$ inputs. \emph{Note:} the code
still implements the earlier no-wait window (the \texttt{service} /
\texttt{latest\_start} fields and \texttt{hard\_windows} flag); migrating it to
the reservation model of~\eqref{eq:reserve} is a pending change.

\end{document}
